% Source: https://artofproblemsolving.com/community/c3h1961913_digits_abcd
% Stable identity: aops-topic-1961913-q01
Since $a\ne0$, $\overline{ab}>b$, so $b\ge2$ (for $b=0$ or $1$ the right side is $0$ or $1$ while the left side exceeds $1$), and $\overline{ab}$ and $b$ have exactly the same prime factors. Writing both sides as products of primes, $\overline{ab}^{10c+5}=b^{10d+5}$ forces $\overline{ab}$ and $b$ to be powers of a common integer: precisely, if $b=t^k$ with $t$ not a perfect power, then $\overline{ab}=t^m$ with $m>k$ and $m(10c+5)=k(10d+5)$.

\emph{The two-digit number $\overline{ab}$ ends in the digit $b$ and is a power of the same base.} Going through $b=2,\dots,9$:
\begin{itemize}
\item $b=2$: two-digit powers of $2$ ending in $2$: $32=2^5$.
\item $b=3$: powers of $3$ ending in $3$ are $3,243,\dots$; none has two digits.
\item $b=4=2^2$: two-digit powers of $2$ ending in $4$: $64=2^6$.
\item $b=5$: $25=5^2$; $b=6$: $16$ and $36$ end in $6$, and $36=6^2$ (while $16$ is not a power of $6$).
\item $b=7,8,9$: two-digit powers of $7$, $2$, $3$ ending in $7,8,9$ do not exist ($49$, $27$, $81$ end in other digits; $8=2^3$ but $2^m$ ends in $8$ only for $m=3,7,\dots$, giving $8,128,\dots$).
\end{itemize}
Now match exponents, remembering that $10c+5$ and $10d+5$ are odd:
\begin{itemize}
\item $b=5$, $\overline{ab}=25$: $2(10c+5)=10d+5$ is impossible (even $=$ odd). Similarly $b=6$, $\overline{ab}=36$ is impossible.
\item $b=2$, $\overline{ab}=32$: $5(10c+5)=10d+5$, i.e. $d=5c+2$; with $1\le c$, $d\le9$ this gives $c=1$, $d=7$.
\item $b=4$, $\overline{ab}=64$: $2^{6(10c+5)}=2^{2(10d+5)}$, i.e. $3(10c+5)=10d+5$, so $d=3c+1$: $c=1,d=4$ or $c=2,d=7$.
\end{itemize}
Therefore $(a,b,c,d)\in\boxed{\{(3,2,1,7),\ (6,4,1,4),\ (6,4,2,7)\}}$, that is, $32^{15}=2^{75}$, $64^{15}=4^{45}$ and $64^{25}=4^{75}$. Checked by computer.
